% BHU PYQ -- Auto-generated & error-corrected
\documentclass[11pt,a4paper]{article}

\usepackage[a4paper,top=2.5cm,bottom=2.5cm,left=2.8cm,right=2.8cm,headheight=14pt]{geometry}
\usepackage{amsmath,amssymb}
\usepackage[T1]{fontenc}
\usepackage[utf8]{inputenc}
\usepackage{lmodern}
\usepackage{microtype}
\usepackage{fancyhdr}
\usepackage{enumitem}
\usepackage{hyperref}
\usepackage{lastpage}
\usepackage{parskip}

\hypersetup{colorlinks=true,urlcolor=black,linkcolor=black,
  pdftitle={STB-401: Sample Surveys and Design of Experiments -- BHU PYQ},pdfauthor={Banaras Hindu University}}

\pagestyle{fancy}\fancyhf{}
\renewcommand{\headrulewidth}{0.4pt}\renewcommand{\footrulewidth}{0.4pt}
\fancyhead[L]{\small   \textit{Banaras Hindu University}}
\fancyhead[R]{\small   \textit{STB-401: Sample Surveys and Design of Experiments}}
\fancyfoot[C]{\small Page  \thepage\ of \pageref{LastPage}}


\newlist{parts}{enumerate}{2}
\setlist[parts,1]{label=(\alph*),leftmargin=*,itemsep=5pt,topsep=3pt}
\setlist[parts,2]{label=(\roman*),leftmargin=*,itemsep=3pt,topsep=2pt}

\newcommand{\pts}[1]{\hfill   \textit{[#1]}}


\begin{document}


\begin{center}
  {\large\bfseries BANARAS HINDU UNIVERSITY}\\[2pt]
  {\normalsize B.Sc. (Hons.) Semester V Examination, 2022-23}\\[6pt]
  \rule{0.8\linewidth}{0.5pt}\\[6pt]
  {\large\bfseries Paper No. STB-401: Sample Surveys and Design of Experiments}\\[4pt]
  {\normalsize Time: 3 Hours\hfill Maximum Marks: 70}
\end{center}
\rule{\linewidth}{0.4pt}
\smallskip



\noindent   \textit{Instructions: Answer   \textbf{five} questions including Question~1 which is compulsory. All questions carry equal marks. Symbols have their usual meanings.}
\bigskip

\medskip


\noindent   \textbf{Question 1 (Compulsory).}\pts{14}


\begin{parts}
  \item The households in a town are to be sampled in order to estimate the average amount of assets per household. The households are stratified into a high-rent stratum and a low-rent stratum. A household in the high-rent stratum is thought to have about nine times as much assets as one in the low-rent stratum, and the stratum standard deviation $S_h$ is expected to be proportional to the square root of the stratum mean. There are 4,000 households in the high-rent stratum and 20,000 in the low-rent stratum. Distribute a sample of 1,000 households between the two strata using optimum allocation. \pts{7}
  \item In a stratification with two strata, $W_1 = 0.8$, $W_2 = 0.2$, $S_1 = 2$ and $S_2 = 4$. Compute the sample sizes $n_1$ and $n_2$ (ignoring the finite population correction) if the standard error of the estimated population mean $ar{y}_{st}$ is to be $0.1$ and the total sample size $n = n_1 + n_2$ is to be minimized. \pts{7}
\end{parts}

\medskip\rule{\linewidth}{0.2pt}

\medskip


\noindent   \textbf{Question 2.}

\begin{parts}
  \item Discuss in detail the different types of errors involved in sample surveys. \pts{7}
  \item Consider a population of 6 units with values $1, 2, 3, 4, 5, 6$. List all possible samples of size 2 (without replacement) from this population. Verify that the sample mean is an unbiased estimator of the population mean. Also calculate the sampling variance of the sample mean and verify that:
  
\begin{parts}
    \item it agrees with the formula for the variance of the sample mean under SRSWOR, and
    \item this variance is less than the variance of the sample mean obtained under SRSWR.
  \end{parts}\pts{7}
\end{parts}

\medskip\rule{\linewidth}{0.2pt}

\medskip


\noindent   \textbf{Question 3.}

\begin{parts}
  \item Discuss the process of determination of sample size for obtaining an estimate with a specified coefficient of variation. \pts{7}
  \item Discuss ratio, regression, and product methods of estimation. Mention their relative importance in sampling theory. \pts{7}
\end{parts}

\medskip\rule{\linewidth}{0.2pt}

\medskip


\noindent   \textbf{Question 4.}

\begin{parts}
  \item Describe the Analysis of Variance (ANOVA) for the fixed-effects model for one-way classified data. Also discuss the use and method of obtaining the critical difference (CD) in this context. \pts{7}
  \item Mention the advantages and disadvantages of the Completely Randomised Design (CRD), Randomised Block Design (RBD), and Latin Square Design (LSD). \pts{7}
\end{parts}

\medskip\rule{\linewidth}{0.2pt}

\medskip


\noindent   \textbf{Question 5.}

\begin{parts}
  \item Explain the procedure of estimating one missing observation in an RBD. \pts{7}
  \item What are the basic principles of design of experiments? Explain each with a suitable example. Also discuss the layout of a Latin Square Design. \pts{7}
\end{parts}

\medskip\rule{\linewidth}{0.2pt}

\medskip


\noindent   \textbf{Question 6.}

\begin{parts}
  \item Describe stratified random sampling. Write its advantages over simple random sampling and give the formula for estimating the population mean. \pts{7}
\end{parts}

\medskip\rule{\linewidth}{0.2pt}

\medskip


\noindent   \textbf{Question 7.}

What are the advantages of:

\begin{parts}
  \item simple random sampling over systematic sampling?
  \item LSD over RBD?
  \item objective sampling over subjective sampling?
  \item sampling over complete enumeration?
\end{parts}\pts{14}

\medskip\rule{\linewidth}{0.2pt}

\medskip


\noindent   \textbf{Question 8.} Write short notes on any   \textbf{two} of the following:\pts{14}

\begin{parts}
  \item Systematic sampling
  \item Factorial experiments
  \item $F$-test and its applications
  \item Cluster sampling and two-stage sampling
\end{parts}

\vfill
\rule{\linewidth}{0.4pt}

\begin{center}\small
  \href{https://bhu-exam.vercel.app/}{Click here to visit our website}\\[4pt]
  \href{https://me-aryan.vercel.app/}{Click here to visit our contributor}\\[4pt]
  \href{https://quashan-me.vercel.app/}{Click here to visit our contributor}
\end{center}

\end{document}